%! AP Statistics — Unit 5, Lesson 5.2 — Homework
\documentclass[10pt]{article}
\usepackage{apstats-article}
\usepackage{apstats-boxes}
\newcommand{\TallMath}[1]{$\displaystyle #1\rule[-1.4em]{0pt}{3.2em}$}

\begin{document}

\pageheader{Unit 5, Lesson 5.2}{Homework}
\namedateperiod

\vspace{0.15in}

For each problem, draw a sketch of the normal curve, shade the appropriate region, and
show all calculations.

\begin{enumerate}

  \item Women's heights are approximately $N(64.0, 2.8)$ inches.
        \begin{enumerate}[label=\alph*.]
          \item Find $P(X < 60)$.
                \vspace{1.0in}
          \item Find $P(62 < X < 68)$.
                \vspace{1.0in}
          \item What height is the 90th percentile for women?
                \vspace{0.9in}
          \item Interpret your answer to (c) in a complete sentence.
                \writelines{2}
        \end{enumerate}

  \item The masses of eggs from a hen are approximately $N(62\text{ g}, 5\text{ g})$.
        \begin{enumerate}[label=\alph*.]
          \item Large eggs are those above 63 g. Find $P(X > 63)$.
                \vspace{1.0in}
          \item Extra-large eggs are the top 15\%. Find the minimum mass for an
                extra-large egg.
                \vspace{0.9in}
        \end{enumerate}

  \item A student claims that the number of text messages sent per day follows a normal
        distribution. A histogram of 200 students' daily text messages shows a
        distribution that is strongly right-skewed with a long tail toward high values.
        \begin{enumerate}[label=\alph*.]
          \item Is the normal model appropriate here? Explain.
                \writelines{2}
          \item What could the student do to make the normal model more appropriate?
                \textit{(Hint: think about sample size and the CLT --- preview for
                Lesson 5.3.)}
                \writelines{2}
        \end{enumerate}

\end{enumerate}

\begin{extensionbox}
\textbf{Extension (optional).}
Cholesterol levels for adults are approximately $N(200, 35)$ mg/dL. The most extreme
4\% of cholesterol levels (2\% in each tail) require medical monitoring. Find the
two boundary values and interpret both in context. (\textit{Hint:} VAR-6.B.2d ---
split the 4\% equally into two tails.)
\end{extensionbox}

\begin{reflectionbox}
\textbf{Preview:} Lesson 5.3 introduces the Central Limit Theorem, which explains why
sampling distributions of means are approximately normal --- even for non-normal
populations --- when samples are large enough.
\end{reflectionbox}

\end{document}
